\section{Estructura preliminar de la tesis}
\label{sec:estructura-tesis}

% Descripción general de los capítulos, y su contenido, que serán
% incluidos en el documento de tesis (según plantilla institucional).
% Nota: esto describe la TESIS completa (el documento final), no el
% protocolo actual -- es un adelanto de cómo se organizará ese documento.

\begin{enumerate}
  \item \textbf{Capítulo 1. Introducción} --
    El primer capítulo corresponde al protocolo de investigación consolidado y enmarca todo el trabajo. Se establece el contexto y la motivación, partiendo del desarrollo de software asistido por IA y del Role Specialization Model (RSM) propuesto en un trabajo previo, y se plantea el problema: la falta de evidencia empírica y cuantitativa sobre su efectividad. A partir de ahí se formula la hipótesis y se definen el objetivo general y los cuatro objetivos específicos, se delimitan el alcance y las limitaciones del estudio, y se describe la metodología propuesta para desarrollar y evaluar el artefacto de investigación, que los capítulos posteriores llevan a cabo.

  \item \textbf{Capítulo 2. Estado del arte} --
    El segundo capítulo desarrolla el objetivo O1 y sienta las bases conceptuales del trabajo. Se revisa la evolución del desarrollo de software asistido por IA, desde los asistentes que actúan como copiloto hasta los agentes de codificación que razonan sobre la intención de la tarea, y en ese panorama se sitúa la orquestación de múltiples agentes basados en modelos de lenguaje. Se analiza el Role Specialization Model (RSM) y los trabajos sobre especialización de roles y heterogeneidad de modelos, identificando qué evidencia existe y qué vacíos deja el estudio exploratorio previo. A partir de esa revisión se seleccionan y justifican las métricas con las que se compararán ambos enfoques: las de corrección funcional (compilabilidad y aprobación de pruebas) y las de estructura arquitectónica (modularidad y complejidad cognitiva), vinculadas con las características de la norma ISO/IEC 25010.

  \item \textbf{Capítulo 3. Diseño e implementación del artefacto} --
    El tercer capítulo aborda el objetivo O2 y constituye el núcleo de la aportación: el diseño y la implementación del artefacto, el modelo de orquestación multiagente basado en el RSM con heterogeneidad completa de modelos. Se parte de los requisitos de diseño derivados del problema y se describe la arquitectura del modelo de orquestación: los tres roles del RSM (arquitecto, analista y especialista), el flujo de coordinación mediante el cual sus contribuciones se integran en una solución única por tarea, y el acceso a modelos de distintas familias o proveedores a través de un enrutador de modelos en la nube. Se detallan los criterios de asignación de los roles a los modelos (el de mayor capacidad en el rol de arquitecto, y modelos de otras dos familias en los roles de analista y especialista) y las decisiones de instrumentación que hacen operable el flujo. Por último, se describe la configuración de la línea base con la que se contrasta el artefacto: un agente único de propósito general que ejecuta todas las fases del desarrollo con el mismo modelo asignado al rol de arquitecto, de modo que el modelo más capaz esté presente en ambas condiciones y la diferencia no pueda atribuirse únicamente a la capacidad del modelo del arquitecto. Finalmente, el capítulo cubre la preparación y validación del aparato experimental: la construcción del banco de tareas (funcionalidades de backend en TypeScript, cada una un caso de uso, con sus suites de prueba) y un piloto que valida el artefacto y el arnés sobre un subconjunto de tareas y calibra el número de repeticiones por tarea (K).

  \item \textbf{Capítulo 4. Evaluación experimental y resultados} --
    El cuarto capítulo cubre los objetivos O3 y O4. En su primera parte se detalla el diseño experimental: las variables dependientes y su operacionalización, y las condiciones que se mantienen constantes para que la comparación sea equitativa. Se describe el experimento confirmatorio, que compara ambos enfoques sobre el banco de tareas construido en el capítulo anterior, con K repeticiones por tarea bajo un diseño pareado. En su segunda parte se reportan los resultados: la ejecución del experimento, el análisis estadístico de cada métrica (pruebas de normalidad, contraste pareado y tamaño del efecto) y la discusión de los hallazgos frente a la hipótesis y a las amenazas a la validez.

  \item \textbf{Capítulo 5. Conclusiones} --
    El quinto capítulo sintetiza las conclusiones de la investigación y responde a la hipótesis a partir de la evidencia obtenida. Se discute en qué medida la combinación de roles especializados con heterogeneidad completa de modelos influye en la corrección funcional y la estructura del código frente a un agente único de propósito general. Se presentan las aportaciones del trabajo en los planos práctico, metodológico y teórico, y se exponen sus limitaciones y las líneas de trabajo futuro que abre.
\end{enumerate}
